\documentclass[11pt, a4paper]{article}
\usepackage{amsmath, amsthm, amssymb}
\usepackage[margin=2.5cm]{geometry}
\usepackage{hyperref}

\newtheorem{theorem}{Theorem}
\newtheorem{corollary}[theorem]{Corollary}
\newtheorem{proposition}[theorem]{Proposition}
\newtheorem*{remark}{Remark}

\title{Discrete Preservation of Radial and Azimuthal Symmetry\\[4pt]
       \large A note on the Petrov--Galerkin Landau discretization}
\author{}
\date{}

\begin{document}
\maketitle

\section*{Setup and notation}

We work with the Petrov--Galerkin discretization of the Landau collision
operator described in the main text (Sec.~3--5).  Trial (basis) and test
functions are indexed by triples $(k,l,m)$ with $k \geq 0$,
$0 \leq l \leq n-1$, $|m| \leq l$:
\[
  \psi_{k,l,m}(p)
    = \mu_{k,l}\, L_k^{l+1/2}(r^2)\, r^l\, Y_{lm}(\hat p)\, e^{-r^2/2},
  \qquad
  \phi_{k,l,m}(p)
    = L_k^{l+1/2}(r^2)\, r^l\, Y_{lm}(\hat p),
\]
where $r = |p|$, $\hat p = p/r$, $\mu_{k,l}$ is the normalisation constant
(Eq.~(43)), and $Y_{lm}$ are real spherical harmonics (Eqs.~(28)--(30)).
We write $i = (k_i, l_i, m_i)$ for a multi-index and use the same letter for
the corresponding flat index when the context is clear.

The coefficient vector $\mathbf{f} \in \mathbb{R}^N$ ($N = n^3$) evolves by
the forward-Euler scheme
\begin{equation}\label{eq:scheme}
  \mathbf{f}^{n+1} = \mathbf{f}^n
    + \tau\, M^{-1}\, Q(\mathbf{f}^n, \mathbf{f}^n),
\end{equation}
where $M_{ij} = \langle \phi_i, \psi_j \rangle$ is the mass matrix (Eq.~(14))
and
\begin{equation}\label{eq:Q}
  [Q(\mathbf{f}, \mathbf{f})]_i
    = \sum_{s,\, t} f_s\, f_t\; \mathcal{Q}_{i,s,t}
\end{equation}
with $\mathcal{Q}_{i,s,t}$ the collision tensor (Eq.~(15); index order:
$i$ = test function, $s = f(p)$, $t = \nabla g(q)$).

Two sparsity rules from Sec.~7 identify when tensor entries vanish
\emph{identically}:

\begin{itemize}
  \item \textbf{Anisotropic rule} (Sec.~7.3, ``andrea''):
    $\mathcal{Q}_{i,s,t} = 0$ unless $l_s + l_t - l_i$ is a non-negative
    even integer satisfying $l_s + l_t - l_i \leq 2\min(l_s, l_t)$.
    Equivalently, $l_i$ must appear in the Clebsch--Gordan decomposition of
    $l_s \otimes l_t$, i.e.\ $|l_s - l_t| \leq l_i \leq l_s + l_t$ with
    $l_s + l_t - l_i$ even.

  \item \textbf{Directional rule} (Sec.~7.2, ``cai''):
    $\mathcal{Q}_{i,s,t} = 0$ unless $|m_i| = |m_s + m_t|$ or
    $|m_i| = |m_s - m_t|$.
\end{itemize}

\noindent We also record one structural fact about the mass matrix.

\begin{proposition}[Block structure of $M$]\label{prop:block}
$M_{ij} = 0$ whenever $(l_i, m_i) \neq (l_j, m_j)$.  In particular $M$,
and therefore $M^{-1}$, are block-diagonal with blocks indexed by $(l, m)$.
\end{proposition}

\begin{proof}
$M_{ij} = \langle \phi_i, \psi_j \rangle$ factorises into a radial integral
times the angular integral $\int_{S^2} Y_{l_i m_i}(\hat p)\, Y_{l_j m_j}(\hat p)
\,d\hat p = \delta_{l_i l_j}\,\delta_{m_i m_j}$, which vanishes unless
$(l_i, m_i) = (l_j, m_j)$.
\end{proof}

\section*{Main results}

\begin{theorem}[Discrete preservation of radial symmetry]\label{thm:radial}
If $f^0_{k,l,m} = 0$ for every $l > 0$, then $f^n_{k,l,m} = 0$ for every
$l > 0$ and every time step $n \geq 0$.
\end{theorem}

\begin{proof}
By induction it suffices to show: if $f_{k,l,m} = 0$ for all $l > 0$ then
$[\mathbf{f}^{n+1} - \mathbf{f}^n]_i = \tau[M^{-1} Q(\mathbf{f},\mathbf{f})]_i
= 0$ for every $i$ with $l_i > 0$.

\textit{Step 1: $[Q(\mathbf{f},\mathbf{f})]_i = 0$ for $l_i > 0$.}
By hypothesis $f_s = 0$ unless $l_s = 0$, and similarly $f_t = 0$ unless
$l_t = 0$, so the only potentially non-zero terms in~\eqref{eq:Q} have
$l_s = l_t = 0$.  For any such pair and any $i$ with $l_i > 0$, the
anisotropic rule requires
\[
  0 \;\leq\; \underbrace{l_s + l_t}_{=\,0} - l_i \;=\; -l_i,
\]
which fails for $l_i > 0$.  Hence $\mathcal{Q}_{i,s,t} = 0$ for every
$(s,t)$ with $l_s = l_t = 0$ and every $i$ with $l_i > 0$, and therefore
$[Q(\mathbf{f},\mathbf{f})]_i = 0$.

\textit{Step 2: $[M^{-1} Q(\mathbf{f},\mathbf{f})]_i = 0$ for $l_i > 0$.}
By Step~1, $Q(\mathbf{f},\mathbf{f})$ has non-zero entries only at $l = 0$
indices.  By Proposition~\ref{prop:block}, $M^{-1}$ is block-diagonal in
$(l,m)$, so row $i$ of $M^{-1}$ is supported on the $(l_i, m_i)$-block.
For $l_i > 0$ that block of $Q(\mathbf{f},\mathbf{f})$ is zero, giving
$[M^{-1} Q(\mathbf{f},\mathbf{f})]_i = 0$.
\end{proof}

\begin{remark}
The zeros in Step~1 are \emph{structural}: each tensor entry
$\mathcal{Q}_{i,s,t}$ with $l_s = l_t = 0$ and $l_i > 0$ vanishes
identically, not by cancellation among non-zero summands.  In exact
arithmetic the symmetry is exact; in floating-point arithmetic it is
preserved to machine precision.
\end{remark}

\begin{corollary}[Discrete preservation of azimuthal symmetry]\label{cor:azimuthal}
If $f^0_{k,l,m} = 0$ for every $m \neq 0$, then $f^n_{k,l,m} = 0$ for
every $m \neq 0$ and every $n \geq 0$.
\end{corollary}

\begin{proof}
The only potentially non-zero terms in~\eqref{eq:Q} have $m_s = m_t = 0$.
The directional rule then requires $|m_i| = |0 \pm 0| = 0$, so
$\mathcal{Q}_{i,s,t} = 0$ for every $|m_i| > 0$.  The block-diagonal
structure of $M^{-1}$ (in $m$) completes the argument.
\end{proof}

\begin{remark}[Physical interpretation and generality]
Both results are discrete manifestations of the rotational equivariance of
the Landau kernel.  The tensor $\mathbb{S}(u)$ (Eq.~(10)) satisfies
$\mathbb{S}(Ru) = R\,\mathbb{S}(u)\,R^\top$ for every $R \in \mathrm{SO}(3)$,
so the weak form commutes with simultaneous rotation of $p$ and $q$.  The
Wigner--Eckart theorem then forces the angular coupling to respect the
Clebsch--Gordan rules, which is precisely the anisotropic selection rule.
For $l_s = l_t = 0$ the Clebsch--Gordan product $0 \otimes 0$ contains only
$l = 0$, closing off all $l > 0$ channels.

Consequently, Theorem~\ref{thm:radial} and Corollary~\ref{cor:azimuthal}
hold for \emph{any} Petrov--Galerkin scheme in the spherical-harmonic--Laguerre
basis applied to any rotation-equivariant bilinear collision operator, not
just the Landau operator with the simplified kernel $\Phi_{\mathrm{simple}}$.
The only requirements are the Clebsch--Gordan selection rule on the collision
tensor and the block-diagonal structure of the mass matrix, both of which
follow from the rotational symmetry of the basis and the kernel.
\end{remark}

\begin{remark}[Invariant subspaces]
More generally, any set $\mathcal{S}$ of values of $l$ that is
\emph{closed under the triangle rule} --- meaning $l_s, l_t \in \mathcal{S}$
implies every $l$ with $|l_s - l_t| \leq l \leq l_s + l_t$ and
$l_s + l_t - l$ even lies in $\mathcal{S}$ --- defines an invariant subspace
of the discrete dynamics.  The set $\{0\}$ (radial symmetry) and the set
$\{0, 1, 2, \ldots\}$ with $m = 0$ only (azimuthal symmetry) are both closed
in this sense.
\end{remark}

\end{document}
